\worldchapter{Würfe und Proben}{rolls}
\label{ch:rolls}

\begin{multicols}{2}

Immer wenn ein Charakter eine Aktion mit ungewissem Ausgang durchführt, wird auf einem entsprechenden Wert gewürfelt, der in der Regel vom Spielleiter festgelegt wird. Das Ergebnis des Wurfes zeigt an, ob und wie gut die Handlung gelungen ist. Dies wird als \textit{Wurf} oder \textit{Probe} bezeichnet.

Alle Würfe werden mit sechsseitigen Würfeln durchgeführt. Die Anzahl der Würfel richtet sich nach dem Wert des Charakters in der jeweiligen Eigenschaft, hinzu kommen eventuell Bonus- oder Schicksalswürfel.

Eine Probe ist also ein Wurf mit einer \textit{Anzahl} von Würfeln, wobei ein \textit{Mindestwurf} erreicht werden muss, der durch einen \textit{Schwierigkeitsgrad} modifiziert werden kann. Im Folgenden werden die Elemente näher erläutert.

Diese Art von Proben wird überall außerhalb des Kampfes verwendet. Im Kampf gelten besondere Regeln für die Bestimmung von Treffern und Verletzungen.

\section{Die Anzahl der Würfel}

Die Anzahl der Würfel für eine Probe entspricht exakt dem Gesamtwert der geforderten Fertigkeit. Dieser Gesamtwert ist die Summe aus dem Basisattribut der Fertigkeit und allen Boni, die durch Charakterschablonen hinzugefügt wurden.

\begin{pexample}[Beispiel]
Ein Charakter mit einem Gesamtwert von 4 in ``Einschüchtern'' (z.B. durch Auffassungsgabe 1 + Schablone 3) würfelt mit 4 sechsseitigen Würfeln.
\end{pexample}

Die Wissenseigenschaften haben einen Wert, der die Anzahl der Würfel bestimmt. Zusätzlich ist aber noch eine zugehörige Fertigkeit angegeben, deren Wert zum Wissenswert addiert wird. Ein Charakter mit der Fertigkeit ``Kommunikation'' 2 und dem Wissen ``Etikette (Kommunikation)'' 3 hat also insgesamt 5 Würfel auf Etikette.

Ist der Gesamtwert 0 oder negativ, kann der Charakter den Wurf nicht ohne weitere Hilfe durchführen, er ist einfach zu schlecht in dieser Fertigkeit. \textit{Bonuswürfel} oder \textit{Schicksalswürfel} können jedoch verwendet werden, auch wenn der Wert negativ ist.

\section{Der Mindestwurf}

Der Mindestwurf eines Charakters ist 5+. Dieser kann durch die Abstammung, weitere Charakterschablonen oder besondere Umstände verändert sein.

\begin{pexample}[Beispiel]
Der Spieler von Hagen möchte mit roher Gewalt eine Tür aufbrechen, die zwischen ihm und der vermuteten Diebesbande steht. Der Spielleiter fordert ihn auf, auf Kraft zu würfeln.

Hagen hat einen Wert von 4, er würfelt also mit 4 Würfeln. Jeder Würfel, der eine 5 oder höher zeigt, ist ein Erfolg. Hagens Spieler würfelt jedoch 4 Erfolge und tritt sofort dem ersten Dieb die Tür ins Gesicht.
\end{pexample}

Der Mindestentwurf wird in der Form ``X+'' angegeben, um anzuzeigen, dass es sich um das Würfelergebnis handelt, das mindestens erreicht werden muss.

\section{Schwierigkeitsgrade}

Der Spielleiter kann den Mindestwurf für besonders leichte oder schwere Proben modifizieren. Bei schweren Proben kann der Modifikator als Wurf +, bei leichten Proben als Wurf - angegeben werden.

Eine +3 Probe bedeutet, dass der Mindestwurf um 3 erhöht wird, in der Regel also auf 8+. Hier kommt ins Spiel, dass alle Würfe \textit{weiter} gewürfelt werden, also \textit{explodierende Würfel} sind.

Übliche Schwierigkeitsgrade sind:

\begin{itemize}
    \item \textbf{-2}: sehr leicht
    \item \textbf{-1}: leicht
    \item \textbf{0}: normal
    \item \textbf{+1}: schwer
    \item \textbf{+3}: sehr schwer
    \item \textbf{+6}: extrem schwer
    \item \textbf{+12}: unmöglich
\end{itemize}

\begin{pexample}[Beispiel]
Zwei Schlösser sind zu knacken, ein einfaches Vorhängeschloss und ein komplexes Zylinderschloss. Hagen hat eine Fertigkeit von 3 in ``Schlösser knacken''. Der Spielleiter verlangt für das Vorhängeschloss eine Probe von -1, für das Zylinderschloss eine Probe von +6.

Hagen würfelt mit 3 Würfeln auf 4+ für das Vorhängeschloss, mit 3 Würfeln auf 11+ für das Zylinderschloss.
\end{pexample}

\section{Explodierende Würfel}

In \trans{book_title} ist es möglich, dass die Mindestwürfe höher als 6+ sind, manchmal sogar deutlich höher. Hier gilt für jeden Wurf die Regel der \textit{explodierenden Würfel}.

Würfel, die nach dem Wurf eine 6 ergeben, dürfen noch einmal geworfen werden. Das Ergebnis wird dann addiert. So kann z.B. auf einen Würfel eine 9+ gewürfelt werden, indem zuerst eine 6 und dann mindestens eine 3 gewürfelt wird. Eine 14+ kann nur erreicht werden, wenn auf einem Würfel zuerst eine 6, dann eine weitere 6 und dann mindestens eine 2 gewürfelt wird.

Da bei einem Wurf mit mehreren Würfeln nicht zwischen den einzelnen Würfeln unterschieden wird, können alle Sechsen gleichzeitig neu gewürfelt werden, wenn dies zur Erzielung eines hohen Mindestwurfs erforderlich ist.

\begin{pexample}[Beispiel]
Hagens Spielleiter verlangt von ihm eine Tapferkeitsprobe +9, da er allein gegen die Räuberbande kämpft. Er muss also mit mindestens einem seiner Tapferkeitswürfel eine 14 würfeln. Glücklicherweise hat er einen Tapferkeitswert von 5, so dass ihm 5 Würfel zur Verfügung stehen.

Im ersten Wurf würfelt er 4,2,6,6,1. Er hat also zwei Sechsen zur Verfügung, die er weiter würfeln kann, um eventuell noch die 14 zu erreichen. Der zweite Wurf (mit den beiden Würfeln) zeigt eine 6 und eine 1.

Jetzt darf er nur noch den verbliebenen Würfel, der eine 6 zeigt, neu würfeln. Da dieser beim zweiten Wurf nur eine 1 zeigt, hilft ihm auch die ausgeprägte Tapferkeit nicht, Hagen erreicht nur eine 13.
\end{pexample}

\section{Kritische Erfolge}

Ähnlich wie bei \textit{kritischen Treffern} im Kampf kann es auch bei anderen Würfen zu kritischen Erfolgen kommen. Erreicht ein Würfel beim Weiterwürfeln mindestens eine 11, so ist dies ein kritischer Erfolg. Dies entspricht einem ``weiter gewürfelten'' \textit{Explodierenden Würfel}, der danach wieder ein Ergebnis von 5+ zeigt. Änderungen am \textit{Mindestwurf} des Charakters werden hier nicht angewendet.

Kritische Würfelerfolge ergeben einen zusätzlichen Erfolg für jedes Mal, wenn nach einem weiteren Würfelwurf eine 5+ erreicht wird. Daraus ergeben sich folgende Grenzen für zusätzliche Erfolge.

\begin{itemize}
    \item \textbf{Wurf 5+}: normaler Erfolg
    \item \textbf{Wurf 11+}: kritischer Erfolg - ergibt einen zusätzlichen Erfolg
    \item \textbf{Wurf 17+}: megakritischer Erfolg - ergibt zwei zusätzliche Erfolge
    \item \textbf{Wurf 23+}: megakritischer Erfolg - ergibt drei zusätzliche Erfolge
    \item \textbf{Wurf 29+}: megakritischer Erfolg - ergibt vier zusätzliche Erfolge
    \item usw.
\end{itemize}

\section{Bonuswürfel}

Ein Charakter kann eine bestimmte Anzahl von Bonuswürfeln haben. Diese werden durch die Schablonen festgelegt (siehe \cref{ch:create}). Bonuswürfel können in beliebiger Anzahl zu den Würfeln eines Wurfes hinzugefügt werden. Dies kann auch geschehen, wenn der eigentliche Wurf bereits fehlgeschlagen ist. So kann ein Bonuswürfel nach dem anderen geopfert werden, um eventuell doch noch einen Erfolg zu erzielen.

Auf diese Weise kann eine Probe auch gewürfelt werden, wenn die Würfelzahl 0 oder weniger beträgt.

Die Bonuswürfel frischen sich bei jeder Rast auf ihr Maximum auf.

\section{Wiederholungswürfe}

Die Anzahl an Wiederholungswürfen eines Charakters ist eine Charaktereigenschaft, die durch Charakterschablonen erlangt werden kann.

Für jeden Wiederholungswurf kann ein kompletter Wurf erneut geworfen werden. Es ist also nicht möglich, einen Wurf mit einem Wert von 0 oder weniger zu bestehen.

Wiederholungswürfe frischen sich auch bei jeder Rast auf ihr Maximum auf.

\section{Schicksalswürfel}

In der Regel erhält der Charakter Schicksalswürfel mit seinem Werdegang, der Spielleiter kann aber auch einzelne Schicksalswürfel für besondere Aktionen oder zu besonderen Anlässen vergeben.

Schicksalswürfel können sowohl als Bonuswürfel als auch für Wiederholungswürfe verwendet werden. Auf einem Schicksalswürfel ist ein Ergebnis von 4+ \textbf{immer} ein Erfolg, unabhängig vom Schwierigkeitsgrad der Probe. Wird der Schicksalswürfel als Wiederholungswurf verwendet, gilt dies für alle Würfel, die in diesem Wurf geworfen werden.

Die Schicksalswürfel müssen immer getrennt von den normalen Würfeln geworfen werden, um festzustellen, ob sie eine 4+ erreicht haben.

\begin{pexample}[Beispiel]
Betrachtet man das vorherige Beispiel mit Hagens Tapferkeitswurf (5 Würfel auf 14+), so würde ihm hier ein Schicksalswurf sehr helfen, da er eigentlich nur eine 4+ würfeln müsste.
\end{pexample}

Schicksalswürfel frischen sich bei jeder Rast auf ihr Maximum auf.

\section{Gruppenwürfe}

Immer wenn es darum geht, dass die Gruppe als Ganzes eine Probe bestehen muss, werden Gruppenwürfe verwendet. Anstatt z.B. von jedem einzelnen Spieler eine Achtsamkeitsprobe zu verlangen, kann der Spielleiter eine Achtsamkeitsprobe von der ganzen Gruppe verlangen. Wenn diese Probe erfolgreich ist, gilt der Effekt für alle Charaktere der Gruppe.

Ein erfolgreicher Gruppenwurf erfordert immer eine bestimmte Anzahl von Erfolgen. Der Spielleiter sagt an, wie viele Erfolge erzielt werden müssen, damit ein Gruppenwurf erfolgreich ist. Dann würfelt jeder Spieler auf das geforderte Attribut oder die geforderte Fertigkeit.

Alle Erfolge der Charaktere werden zusammengezählt. Ist die erforderliche Anzahl von Erfolgen erreicht, ist die Probe gelungen.

Bei einem Gruppenwurf kann jeder Teilnehmer wie gewohnt Bonus-, Schicksals- und Wiederholungswürfel einsetzen. Es ist auch möglich, den Mindestwurf entsprechend der Schwere der Probe zu ändern.

\section{Verborgenheit}

Ein besonderer Wurf ist der Wurf auf die Verborgenheit. Dieser kommt zum Tragen, wenn eine Person einen Charakter beobachtet und nach bestimmten auffälligen Ausrüstungsgegenständen sucht. Jeder Ausrüstungsgegenstand hat einen Verborgenheitswert. Je höher dieser Wert ist, desto auffälliger ist dieser Gegenstand zu sehen oder zu erkennen.

Beobachtet nun eine Person einen Charakter oder die ganze Gruppe, so gilt für den Verborgenheitswurf der Gegenstand mit dem \textit{höchsten} Verborgenheitswert.

Der Beobachter würfelt mit einer Anzahl von Würfeln, die seinem Wahrnehmungswert \textit{plus} dem ermittelten Verborgenheitswert entspricht. Bei Erfolg kann der Beobachter ein auffälliges Objekt ausmachen.
\end{multicols}
